A decision tree is essentially a flowchart for making predictions. It works by sorting the data into smaller and smaller groups. In decision tree terminology those groups are called leaves. They are created to be as "pure" as possible using a functions that can quantify said purity.\\
Once the tree is built every leaf node is assigned a value used to predict the target variable of new, unseen data.\\
For Regression: The leaf is assigned the Average (Mean) of the values in the leaf.\\
For Classification: The leaf is assigned the Majority Class. In other words the class that is most common in the leaf. This decision is made utilizing certain principles for tie-breaking.\\
The project, in essence, revolves around a regression problem, that of predicting the policy risk of new customers, hence why the classification method couldn't hold for our scope.


\subsection{Structure}
The structure of this project's Decision Tree implementation follows a classic object oriented design. It is constructed with nodes which function as either:\\
\begin{itemize}
\item Leaf - which holds the final prediction.\\
\item Decision node - which holds the splitting rule (the feature and threshold) as well as pointers to it's left and right children. It represents every group leading up to the leaf; it's children which stand for the smaller subgroups of the current one.\\
\end{itemize}
The DecisionTree class manages the logic of how the tree is built and how it makes decisions. It stores the final structure made of nodes, which is similar to a binary tree, where each node points to either nothing or it's two children nodes.\\

\subsection{Finding Splits}
The most important part of the construction of a decision tree is finding a split. This is the most computationally heavy part. The Decision Tree:
\begin{enumerate}
    \item[1.] Loops through every feature
    \item[2.] Sorts the data by feature
    \item[3.] Tests every possible midpoint between values as a potential "threshold."
    \item[4.] Calculates the Impurity Decrease for each using the selected criterion. It remembers the one that reduced the disorder the most.
\end{enumerate}
Depending on the parameters the tree is built from this in one of two ways.
\subsubsection{Recursive}
The standard way of building a tree. The algorithm starts at the Root. It splits the Root into a Left and Right child. It then moves to the Left child and tries to split it. If it can, it moves into that child's Left child and repeats the process.\\
As consequences of that It dives as deep as possible into the leftmost branch until it hits a stopping condition. Only then does it backtrack to fill in the right-hand side branches.
\subsubsection{Priority Queue}
This is the method used when the tree is limited by the total number of leaves. The recursive approach would not work with such limitation.
Instead of diving deep into one branch, the algorithm looks at all current leaf nodes across the entire tree. The priority queue (heap) stores potential "Impurity Decrease" for every leaf.\\
The algorithm always removes from the queue the node that offers the highest impurity decrease, regardless of where it is in the tree or how deep it is.

\subsection{Stopping Parameters}
Why, what problems they solve\\
\\
Maybe an enumerate of stopping params with simple description of how they work, if that fits, if not use subsubsection for each stopping parameter
\begin{enumerate}
    \item[a.] \textbf{Max Depth} \\
        This is the description?
    \item[b.] \textbf{Minimum Samples to Split} -
    \item[c.] \textbf{Minimum Samples at Leaf} - 
    \item[d.] \textbf{Maximum number of Leaves} -
    \item[e.] \textbf{Minimum Impurity Decrease} -  
\end{enumerate}
